\shead{Finding the percentage one number is of another}
\begin{knowbox}\begin{itemize}[leftmargin=1.4em,itemsep=3pt]
\item To find what percentage $A$ is of $B$: \ divide $A$ by $B$, then write the quotient as a percentage (multiply by $100$ and write $\%$).\par
Example: $\ 9$ is what percentage of $12$? \quad $9:12=0.75=75\%$.
\item Read the question carefully: ``$A$ compared with $B$'' means $B$ is the whole ($100\%$) and we divide by $B$.
\item \textbf{Above or below plan:} if the plan is $100\%$, then $\ \text{actual}:\text{plan}$ tells what percentage of the plan was done; the part above $100\%$ is the percentage \emph{exceeding} the plan.
\item \textbf{Percentage increase or decrease:} $\ (\text{change}):(\text{old amount})$, written as a percentage.
\end{itemize}
\end{knowbox}
\gap
\shead{Finding a percentage of a number}
\begin{knowbox}\begin{itemize}[leftmargin=1.4em,itemsep=3pt]
\item To find $a\%$ of $M$, \ calculate $\ M\times a:100\ $ or $\ M:100\times a$.\par
Example: $\ 15\%$ of $60$ kg is $\ 60\times15:100=9$ (kg).
\item $M:100$ is $1\%$ of $M$; multiplying by $a$ gives $a\%$ of $M$. \ This ``go to $1\%$ first'' idea also helps to work backwards: if $a\%$ of a number is $N$, then $1\%$ of it is $N:a$ and the number is $N:a\times100$.
\end{itemize}
\end{knowbox}
\gap
\shead{Percentages in everyday life}
\begin{knowbox}\begin{itemize}[leftmargin=1.4em,itemsep=3pt]
\item \textbf{Discount:} a price cut of $a\%$ means we pay $(100-a)\%$ of the old price.
\item \textbf{Profit:} profit $=$ selling price $-$ cost price. \ ``Profit of $20\%$'' usually means $20\%$ of the cost price, so the selling price is $120\%$ of the cost price.
\item \textbf{Savings interest:} with an interest rate of $0.5\%$ a month, each month the bank adds $0.5\%$ of the money in the account.
\item \textbf{Salt water:} ``$5\%$ salt'' means the salt weighs $5\%$ of the whole salt water.
\end{itemize}
\note{Always ask: ``percentage of \emph{what}?'' -- that quantity is the $100\%$.}
\end{knowbox}
\gap
\begin{exbox}{\ExWord\ 1}
Find the percentage: \ \pa{a} of $0.6$ compared with $2.4$; \ \pa{b} of $7$ compared with $4$.
\guide
\pa{a} $0.6:2.4=0.25=25\%$.\par
\pa{b} $7:4=1.75=175\%$.
\dapso{\pa{a} $25\%$;\ \ \pa{b} $175\%$.}
\end{exbox}
\gap
\begin{exbox}{\ExWord\ 2}
Class 5B has $18$ boys and $22$ girls. What percentage of the class are boys? What percentage are girls?
\guide
Pupils in the class: $18+22=40$.\par
Boys compared with the whole class: $18:40=0.45=45\%$.\par
Girls compared with the whole class:\par
\emph{Method 1.} $22:40=0.55=55\%$. \qquad \emph{Method 2.} $100\%-45\%=55\%$.
\dapso{boys $45\%$, girls $55\%$.}
\end{exbox}
\gap
\begin{exbox}{\ExWord\ 3}
Find $35\%$ of $240$.
\guide
\emph{Method 1.} $240:100\times35=2.4\times35=84$.\par
\emph{Method 2.} $240\times35:100=8\,400:100=84$.
\dapso{$84$.}
\end{exbox}
\gap
\begin{exbox}{\ExWord\ 4}
A youth group planned to plant $250$ trees, but actually planted $290$ trees. What percentage of the plan did they reach, and by what percentage did they exceed the plan?
\guide
Actual compared with plan: $290:250=1.16=116\%$.\par
They exceeded the plan by: $116\%-100\%=16\%$.
\dapso{$116\%$ of the plan; exceeded by $16\%$.}
\end{exbox}
\gap
\begin{exbox}{\ExWord\ 5}
A school bag costs $180\,000$ dong. During a sale its price is cut by $15\%$. How much does the bag cost in the sale?
\guide
\emph{Method 1.} The price cut: $180\,000\times15:100=27\,000$ (dong).\par
Sale price: $180\,000-27\,000=153\,000$ (dong).\par
\emph{Method 2.} We pay $100\%-15\%=85\%$ of the price: $\ 180\,000\times85:100=153\,000$ (dong).
\dapso{$153\,000$ dong.}
\end{exbox}
\dividerbar{Practice}
\begin{probbox}{1}{\relax}
Find the percentage of the first number compared with the second.
\begin{cols3}
\pa{a} $18$ and $24$ & \pa{b} $4.5$ and $12.5$ & \pa{c} $0.9$ and $3.6$\\
\pa{d} $27$ and $15$ & \pa{e} $3.4$ and $8$ & \pa{g} $\dfrac14$ and $5$
\end{cols3}
\end{probbox}
\probsol{
\pa{a} $18:24=0.75=75\%$. \quad \pa{b} $4.5:12.5=0.36=36\%$. \quad \pa{c} $0.9:3.6=0.25=25\%$.\par
\pa{d} $27:15=1.8=180\%$. \quad \pa{e} $3.4:8=0.425=42.5\%$. \quad \pa{g} $\dfrac14:5=0.05=5\%$.
\dapso{$75\%$; $36\%$; $25\%$; $180\%$; $42.5\%$; $5\%$.}}
\gap
\begin{probbox}{2}{\relax}
Calculate.
\begin{cols3}
\pa{a} $15\%$ of $60$ kg & \pa{b} $8\%$ of $450$ m & \pa{c} $25\%$ of $7.2$ l\\
\pa{d} $40\%$ of $3.5$ tonnes & \pa{e} $120\%$ of $85$ cm & \pa{g} $12.5\%$ of $96$ m$^2$
\end{cols3}
\end{probbox}
\probsol{
\pa{a} $60\times15:100=9$ (kg). \quad \pa{b} $450\times8:100=36$ (m). \quad \pa{c} $7.2\times25:100=1.8$ (l).\par
\pa{d} $3.5\times40:100=1.4$ (tonnes). \quad \pa{e} $85\times120:100=102$ (cm).\par
\pa{g} $96\times12.5:100=12$ (m$^2$).
\dapso{$9$ kg; $36$ m; $1.8$ l; $1.4$ tonnes; $102$ cm; $12$ m$^2$.}}
\gap
\begin{probbox}{3}{\relax}
A bag holds $48$ marbles: $30$ are red and the rest are green. \ \pa{a} What percentage of the marbles are red? \ \pa{b} What is the percentage of green marbles compared with red marbles?
\end{probbox}
\probsol{
\pa{a} $30:48=0.625=62.5\%$.\par
\pa{b} Green marbles: $48-30=18$. \ $18:30=0.6=60\%$.
\dapso{\pa{a} $62.5\%$;\ \ \pa{b} $60\%$.}}
\gap
\begin{probbox}{4}{\relax}
In a maths test, $34$ of the $40$ pupils of class 5C got a good mark. What percentage of the class got a good mark?
\end{probbox}
\probsol{
$34:40=0.85=85\%$.
\dapso{$85\%$.}}
\gap
\begin{probbox}{5}{\relax}
A bicycle workshop planned to make $400$ bicycles in a month, but it made $460$. What percentage of the plan did it reach? By what percentage did it exceed the plan?
\end{probbox}
\probsol{
$460:400=1.15=115\%$.\par
Exceeded by: $115\%-100\%=15\%$.
\dapso{$115\%$; exceeded by $15\%$.}}
\gap
\begin{probbox}{6}{\relax}
A village planned to plant $1\,250$ trees this year. By the end of March it had planted $1\,000$ trees. What percentage of the plan has been done, and what percentage is still to do?
\end{probbox}
\probsol{
Done: $1\,000:1\,250=0.8=80\%$ of the plan.\par
Still to do: $100\%-80\%=20\%$ of the plan.
\dapso{$80\%$ done, $20\%$ still to do.}}
\gap
\begin{probbox}{7}{\relax}
Aunt Lan earns $12$ million dong a month and saves $2.7$ million dong of it. What percentage of her pay does she save?
\end{probbox}
\probsol{
$2.7:12=0.225=22.5\%$.
\dapso{$22.5\%$.}}
\gap
\begin{probbox}{8}{\relax}
Last year Minh weighed $38$ kg. This year his weight has gone up by $10\%$ of last year's weight. How much does Minh weigh now?
\end{probbox}
\probsol{
Increase: $38\times10:100=3.8$ (kg).\par
Weight now: $38+3.8=41.8$ (kg).
\dapso{$41.8$ kg.}}
\gap
\begin{probbox}{9}{\relax}
A rectangular plot of land is $50$ m long; its width is $60\%$ of its length. Vegetables are grown on $45\%$ of its area, flowers on $35\%$, and the rest is paths. Find the area of the paths.
\end{probbox}
\probsol{
Width: $50\times60:100=30$ (m). \ Area: $50\times30=1\,500$ (m$^2$).\par
Paths: $100\%-45\%-35\%=20\%$ of the area.\par
Area of the paths: $1\,500\times20:100=300$ (m$^2$).
\dapso{$300$ m$^2$.}}
\gap
\begin{probbox}{10}{\relax}
A shop sold $1\,200$ kg of rice, of which $35\%$ was sticky rice. How many kilograms of ordinary rice did it sell?
\end{probbox}
\probsol{
Ordinary rice: $100\%-35\%=65\%$ of the rice.\par
$1\,200\times65:100=780$ (kg).
\dapso{$780$ kg.}}
\gap
\begin{probbox}{11}{\relax}
A school has $640$ pupils, and $55\%$ of them are girls. How many boys are there?
\end{probbox}
\probsol{
Girls: $640\times55:100=352$.\par
Boys: $640-352=288$.
\dapso{$288$ boys.}}
\gap
\begin{probbox}{12}{\relax}
A shopkeeper buys a shirt for $150\,000$ dong and sells it for $180\,000$ dong. Her profit is what percentage of the cost price?
\end{probbox}
\probsol{
Profit: $180\,000-150\,000=30\,000$ (dong).\par
$30\,000:150\,000=0.2=20\%$.
\dapso{$20\%$ of the cost price.}}
\gap
\begin{probbox}{13}{\relax}
A storybook costs $60\,000$ dong. At the book fair it is sold at a $10\%$ discount. How much does it cost at the fair?
\end{probbox}
\probsol{
We pay $100\%-10\%=90\%$ of the price.\par
$60\,000\times90:100=54\,000$ (dong).
\dapso{$54\,000$ dong.}}
\gap
\begin{probbox}{14}{\relax}
Mrs Hạnh puts $20\,000\,000$ dong into a savings account with interest of $0.6\%$ a month. After one month, how much interest does she get, and how much money is in the account?
\end{probbox}
\probsol{
Interest: $20\,000\,000\times0.6:100=120\,000$ (dong).\par
Money in the account: $20\,000\,000+120\,000=20\,120\,000$ (dong).
\dapso{interest $120\,000$ dong; total $20\,120\,000$ dong.}}
\gap
\begin{probbox}{15}{\relax}
Seawater contains $3.5\%$ salt. How many kilograms of salt are there in $40$ kg of seawater?
\end{probbox}
\probsol{
$40\times3.5:100=1.4$ (kg).
\dapso{$1.4$ kg.}}
\gap
\begin{probbox}{16}{\relax}
The price of a notebook went up from $8\,000$ dong to $9\,000$ dong. By what percentage did the price go up?
\end{probbox}
\probsol{
Increase: $9\,000-8\,000=1\,000$ (dong).\par
Compared with the old price: $1\,000:8\,000=0.125=12.5\%$.
\dapso{$12.5\%$.}}
\gap
\begin{probbox}{17}{\relax}
The price of a pair of trainers went down from $250\,000$ dong to $200\,000$ dong. By what percentage did the price go down?
\end{probbox}
\probsol{
Decrease: $250\,000-200\,000=50\,000$ (dong).\par
Compared with the old price: $50\,000:250\,000=0.2=20\%$.
\dapso{$20\%$.}}
\gap
\begin{probbox}{18}{\relax}
A football team played $25$ matches: it won $16$, drew $4$ and lost the rest. What percentage of the matches did it win, draw and lose?
\end{probbox}
\probsol{
Lost: $25-16-4=5$ (matches).\par
Won: $16:25=0.64=64\%$; \ drew: $4:25=0.16=16\%$; \ lost: $5:25=0.2=20\%$.
\dapso{won $64\%$, drew $16\%$, lost $20\%$.}}
\gap
\begin{probbox}{19}{\relax}
Hoa is reading a $240$-page book. On the first day she reads $30\%$ of the book, on the second day $45\%$. How many pages are left?
\end{probbox}
\probsol{
Pages left: $100\%-30\%-45\%=25\%$ of the book.\par
$240\times25:100=60$ (pages).
\dapso{$60$ pages.}}
\gap
\begin{probbox}{20}{\relax}
A farmer planted $450$ young trees; $18$ of them died. What percentage of the trees survived?
\end{probbox}
\probsol{
Trees that survived: $450-18=432$.\par
$432:450=0.96=96\%$.
\dapso{$96\%$.}}
\gap
\begin{probbox}{21}{\relax}
Out of $500$ bean seeds, $92\%$ sprouted. How many seeds did not sprout?
\end{probbox}
\probsol{
Seeds that sprouted: $500\times92:100=460$.\par
Seeds that did not sprout: $500-460=40$.
\dapso{$40$ seeds.}}
\gap
\begin{probbox}{22}{\relax}
Class 5A has $35$ pupils; $40\%$ of them got an excellent report. How many pupils did not get an excellent report?
\end{probbox}
\probsol{
Excellent: $35\times40:100=14$ (pupils).\par
Not excellent: $35-14=21$ (pupils).
\dapso{$21$ pupils.}}
\gap
\begin{probbox}{23}{\relax}
A basket holds $80$ apples and pears; the pears make up $40\%$. Find the percentage of apples compared with pears.
\end{probbox}
\probsol{
Pears: $80\times40:100=32$. \ Apples: $80-32=48$.\par
Apples compared with pears: $48:32=1.5=150\%$.
\dapso{$150\%$.}}
\gap
\begin{probbox}{24}{\relax}
For a tree-planting day, $32$ of the $40$ pupils of class 5A and $27$ of the $30$ pupils of class 5B took part. Which class took part at the higher rate?
\end{probbox}
\probsol{
Class 5A: $32:40=0.8=80\%$.\par
Class 5B: $27:30=0.9=90\%$.\par
$90\%>80\%$, so class 5B took part at the higher rate.
\dapso{class 5B ($90\%$ against $80\%$).}}
\gap
\begin{probbox}{25}{\relax}
Last year a factory made $2\,400$ products. This year it made $15\%$ more. How many products did it make this year?
\end{probbox}
\probsol{
This year it made $100\%+15\%=115\%$ of last year's number.\par
$2\,400\times115:100=2\,760$ (products).
\dapso{$2\,760$ products.}}
\gap
\begin{probbox}{26}{\relax}
A museum ticket costs $40\,000$ dong. Children get a $25\%$ discount. How much does a family of $2$ adults and $3$ children pay?
\end{probbox}
\probsol{
A child's ticket: $40\,000\times(100-25):100=30\,000$ (dong).\par
The family pays: $40\,000\times2+30\,000\times3=80\,000+90\,000=170\,000$ (dong).
\dapso{$170\,000$ dong.}}
\gap
\begin{probbox}{27}{\relax}
A rectangular garden is $25$ m long and $16$ m wide. The length is increased by $20\%$ and the width stays the same. Find the new area. By what percentage did the area increase?
\end{probbox}
\probsol{
New length: $25\times120:100=30$ (m). \ New area: $30\times16=480$ (m$^2$).\par
Old area: $25\times16=400$ (m$^2$). \ Increase: $480-400=80$ (m$^2$).\par
$80:400=0.2=20\%$.
\dapso{$480$ m$^2$; the area increased by $20\%$.}}
\gap
\begin{probbox}{28}{\relax}
Fresh milk contains $4\%$ fat. How many grams of fat are there in $3$ kg of fresh milk?
\end{probbox}
\probsol{
$3$ kg $=3\,000$ g. \ Fat: $3\,000\times4:100=120$ (g).
\dapso{$120$ g.}}
\gap
\begin{probbox}{29}{\relax}
A lorry has to travel $240$ km in three days. On the first day it covers $35\%$ of the distance, on the second day $40\%$. How many kilometres are left for the third day?
\end{probbox}
\probsol{
Third day: $100\%-35\%-40\%=25\%$ of the distance.\par
$240\times25:100=60$ (km).
\dapso{$60$ km.}}
\gap
\begin{probbox}{30}{\relax}
A class fund has $1\,500\,000$ dong. The class spends $40\%$ of it on books, then $30\%$ of the remaining money on pens. How much money is left?
\end{probbox}
\probsol{
Books: $1\,500\,000\times40:100=600\,000$ (dong).\par
Remaining: $1\,500\,000-600\,000=900\,000$ (dong).\par
Pens: $900\,000\times30:100=270\,000$ (dong).\par
Left: $900\,000-270\,000=630\,000$ (dong).
\dapso{$630\,000$ dong.}}
\dividerbar{Challenge -- Extension}
\begin{probbox}{31}{\stars{1}}
$25\%$ of a number is $18$. Find the number.
\end{probbox}
\probsol{
$1\%$ of the number: $18:25=0.72$.\par
The number ($100\%$): $0.72\times100=72$.
\note{Or: $25\%=\dfrac14$, so the number is $18\times4=72$.}
\dapso{$72$.}}
\gap
\begin{probbox}{32}{\stars{1}}
A factory made $1\,380$ products and so exceeded its plan by $15\%$. How many products did the plan ask for?
\end{probbox}
\probsol{
The factory made $100\%+15\%=115\%$ of the plan.\par
$1\%$ of the plan: $1\,380:115=12$ (products).\par
The plan: $12\times100=1\,200$ (products).
\dapso{$1\,200$ products.}}
\gap
\begin{probbox}{33}{\stars{1}}
A shop sells a fan for $240\,000$ dong and makes a profit of $20\%$ of the cost price. What was the cost price?
\end{probbox}
\probsol{
Selling price: $100\%+20\%=120\%$ of the cost price.\par
Cost price: $240\,000:120\times100=200\,000$ (dong).
\dapso{$200\,000$ dong.}}
\gap
\begin{probbox}{34}{\stars{1}}
After a $25\%$ discount, Mum paid $360\,000$ dong for a jacket. What was the price before the discount?
\end{probbox}
\probsol{
Mum paid $100\%-25\%=75\%$ of the price.\par
Price before the discount: $360\,000:75\times100=480\,000$ (dong).
\dapso{$480\,000$ dong.}}
\gap
\begin{probbox}{35}{\stars{1}}
A shop's profit on a toy is $20\%$ of the \emph{selling} price. What percentage of the \emph{cost} price is the profit?
\end{probbox}
\probsol{
Take the selling price as $100\%$: the profit is $20\%$ and the cost price is $100\%-20\%=80\%$ of the selling price.\par
Profit compared with cost price: $20:80=0.25=25\%$.
\dapso{$25\%$.}}
\gap
\begin{probbox}{36}{\stars{1}}
A shop buys goods at $80\%$ of the listed price and sells them at $5\%$ off the listed price. Its profit is what percentage of the cost price?
\end{probbox}
\probsol{
Take the listed price as $100\%$.\par
Cost price: $80\%$; \ selling price: $100\%-5\%=95\%$; \ profit: $95\%-80\%=15\%$ of the listed price.\par
Profit compared with cost price: $15:80=0.1875=18.75\%$.
\dapso{$18.75\%$.}}
\gap
\begin{probbox}{37}{\stars{1}}
Mr Dũng saves $10\,000\,000$ dong at an interest of $0.5\%$ a month. The interest of the first month is added to the account and also earns interest in the second month. How much money is in the account after two months?
\end{probbox}
\probsol{
After month $1$: interest $10\,000\,000\times0.5:100=50\,000$ (dong); \ total $10\,050\,000$ dong.\par
After month $2$: interest $10\,050\,000\times0.5:100=50\,250$ (dong); \ total $10\,050\,000+50\,250=10\,100\,250$ (dong).
\dapso{$10\,100\,250$ dong.}}
\gap
\begin{probbox}{38}{\stars{1}}
A jar holds $300$ g of salt water containing $6\%$ salt. How many grams of pure water must be added to get salt water containing $4\%$ salt?
\end{probbox}
\probsol{
Salt: $300\times6:100=18$ (g); it does not change when water is added.\par
New salt water: $18$ g is $4\%$ of it, so it weighs $18:4\times100=450$ (g).\par
Water added: $450-300=150$ (g).
\dapso{$150$ g.}}
\gap
\begin{probbox}{39}{\stars{1}}
Mix $200$ g of salt water with $10\%$ salt and $300$ g of salt water with $5\%$ salt. What percentage of the mixture is salt?
\end{probbox}
\probsol{
Salt in the first jar: $200\times10:100=20$ (g); \ in the second: $300\times5:100=15$ (g).\par
Mixture: $200+300=500$ (g), with $20+15=35$ (g) of salt.\par
$35:500=0.07=7\%$.
\dapso{$7\%$.}}
\gap
\begin{probbox}{40}{\stars{1}}
In a class, the boys make up $40\%$ of the pupils. After $6$ girls move to another class, the boys make up $50\%$ of the class. How many pupils did the class have at first?
\end{probbox}
\probsol{
The number of boys does not change.\par
At first: girls $60\%$, boys $40\%$, so girls $=\dfrac{60}{40}=\dfrac32$ of the boys.\par
Afterwards: girls $=$ boys ($50\%$ each).\par
The $6$ girls who left are $\dfrac32-1=\dfrac12$ of the boys, so there are $6\times2=12$ boys.\par
At first the class had $12:40\times100=30$ (pupils).
\dapso{$30$ pupils.}}
\gap
\begin{probbox}{41}{\stars{1}}
The price of a bicycle was raised by $25\%$. By what percentage must the new price be lowered to get back to the old price?
\end{probbox}
\probsol{
Take the old price as $4$ parts; $25\%=\dfrac14$, so the new price is $5$ parts.\par
To get back we must take off $1$ part of the new price: $1:5=0.2=20\%$.
\dapso{$20\%$.}}
\gap
\begin{probbox}{42}{\stars{1}}
Shop A sells a rice cooker for $200\,000$ dong with $15\%$ off. Shop B sells the same rice cooker for $180\,000$ dong with $5\%$ off. Which shop is cheaper, and by how much?
\end{probbox}
\probsol{
Shop A: $200\,000\times85:100=170\,000$ (dong).\par
Shop B: $180\,000\times95:100=171\,000$ (dong).\par
Shop A is cheaper by $171\,000-170\,000=1\,000$ (dong).
\dapso{shop A, cheaper by $1\,000$ dong.}}
\gap
\begin{probbox}{43}{\stars{2}}
A trader sells two vases for $240\,000$ dong each. On one he makes a profit of $20\%$ of its cost; on the other he loses $20\%$ of its cost. Overall, does he gain or lose, and how much?
\end{probbox}
\probsol{
First vase: $240\,000$ is $120\%$ of its cost, so the cost is $240\,000:120\times100=200\,000$ (dong).\par
Second vase: $240\,000$ is $100\%-20\%=80\%$ of its cost, so the cost is $240\,000:80\times100=300\,000$ (dong).\par
Total cost: $500\,000$ dong; \ money received: $480\,000$ dong.\par
He loses $500\,000-480\,000=20\,000$ (dong).
\dapso{he loses $20\,000$ dong.}}
\gap
\begin{probbox}{44}{\stars{1}}
A theatre lowered its ticket price by $20\%$, and the number of people coming went up by $30\%$. By what percentage did the theatre's ticket money change?
\end{probbox}
\probsol{
New price: $80\%$ of the old price; new number of people: $130\%$ of the old number.\par
New ticket money: $80\%\times130\%=0.8\times1.3=1.04=104\%$ of the old ticket money.\par
It went up by $104\%-100\%=4\%$.
\dapso{it went up by $4\%$.}}
\gap
\begin{probbox}{45}{\stars{1}}
The side of a square is increased by $10\%$. By what percentage does its area increase?
\end{probbox}
\probsol{
New side: $110\%$ of the old side.\par
New area: $110\%\times110\%=1.1\times1.1=1.21=121\%$ of the old area.\par
The area increases by $21\%$.
\dapso{$21\%$.}}
\gap
\begin{probbox}{46}{\stars{1}}
A village has $5\,000$ people. Each year the population grows by $2\%$ of the previous year's number. How many people will live in the village after $2$ years?
\end{probbox}
\probsol{
After $1$ year: $5\,000+5\,000\times2:100=5\,000+100=5\,100$ (people).\par
After $2$ years: $5\,100+5\,100\times2:100=5\,100+102=5\,202$ (people).
\dapso{$5\,202$ people.}}
\gap
\begin{probbox}{47}{\stars{2}}
A $500$ g piece of bronze contains $80\%$ copper; the rest is tin. How many grams of copper must be added so that copper makes up $90\%$?
\end{probbox}
\probsol{
Tin: $100\%-80\%=20\%$, \ i.e.\par $500\times20:100=100$ (g). Adding copper does not change the tin.\par
In the new piece, tin is $100\%-90\%=10\%$, so the new piece weighs $100:10\times100=1\,000$ (g).\par
Copper added: $1\,000-500=500$ (g).
\dapso{$500$ g.}}
\gap
\begin{probbox}{48}{\stars{1}}
A stationery shop offers: ``Buy $4$ notebooks, get $1$ free.'' Buying in this way, by what percentage is the price of each notebook reduced?
\end{probbox}
\probsol{
You pay for $4$ notebooks and get $5$. \ The money of $1$ notebook is saved out of the price of $5$ notebooks.\par
$1:5=0.2=20\%$.
\dapso{$20\%$.}}
\gap
\begin{probbox}{49}{\stars{1}}
A saleswoman gets a basic pay of $6\,000\,000$ dong a month plus $5\%$ of the value of the goods she sells. Last month she received $8\,500\,000$ dong. What was the value of the goods she sold?
\end{probbox}
\probsol{
Her extra pay: $8\,500\,000-6\,000\,000=2\,500\,000$ (dong), which is $5\%$ of the goods sold.\par
Value of the goods: $2\,500\,000:5\times100=50\,000\,000$ (dong).
\dapso{$50\,000\,000$ dong.}}
\gap
\begin{probbox}{50}{\stars{2}}
In a maths contest, $60\%$ of the pupils come from the city and the rest from the countryside. $25\%$ of the city pupils and $40\%$ of the countryside pupils win a prize. What percentage of all the pupils win a prize?
\end{probbox}
\probsol{
Countryside pupils: $100\%-60\%=40\%$ of all the pupils.\par
City prize winners: $60\%\times25\%=0.6\times0.25=0.15=15\%$ of all the pupils.\par
Countryside prize winners: $40\%\times40\%=0.4\times0.4=0.16=16\%$ of all the pupils.\par
All prize winners: $15\%+16\%=31\%$.
\dapso{$31\%$.}}
